%==================================================================

\begin{exm}\label{exm:forms}
Пусть $X$ --- гладкое многообразие и $k\geq1$,
$\Shv{F}$ --- предпучок точных ($d$-точных) $k$-форм на $X$,
$\Shv{H}=\Omega^{k}_X$ --- пучок всех $k$-форм. Тогда
$\Shv{F}^{+}$ --- пучок форм, локально точных. По классической лемме
Пуанкаре локально любая замкнутая форма точна, и обратно: точная форма
всегда замкнута ($d^2=0$). Отсюда
\begin{equation}\label{eq:forms}
 \Shv{F}^{+}=\{\text{замкнутые }k\text{-формы}\},\qquad
 \Shv{F}^{\#}\simeq\Shv{F}^{+},
\end{equation}
то есть <<прокачивание>> предпучка точных форм даёт пучок замкнутых
форм: глобально точных форм может быть мало (группы когомологий!),
но локально точных --- столько, сколько хочется.
\end{exm}

\begin{bookbox}
\booksrc{Годеман, гл.~II, \S\,2, с.~156: при доказательстве точности
 используется классическая теорема Пуанкаре (локально любая замкнутая
 форма точна); п.~2.8.1, с.~159 --- примеры пучков $\Omega^p$
 дифференциальных форм и пучков ростков как $\mathcal{O}$-модулей.
 Манин, \S\,2.1.2, с.~110 --- предпучки модулей над предпучком колец.
}
\end{bookbox}

%==================================================================
